%% kapa: exam=1 round=29 year=2018 period=1 form=A subject=civil_law subject_ko=민법 seq=3 number=3
%% source: Q-Net 감정평가사 기출문제 (cst00309) artlSeq=5202808 file=extracted/1008/5202808/2018년 제29회 감정평가사 1차시험 기출문제/2018년 제29회 감정평가사 1차시험 기출문제/1교시 A형.pdf
%% conversion: kordoc 4.24.51 -> tools/gen1.py (auto); review=proofread-sonnet; flags=none
\begin{question}{3}
甲은 취소할 수 없는 법률행위의 범위를 정함이 없이 성년후견개시심판을 받았다. 그 후 甲은 법정대리인 乙의 동의서를 위조하는 방법으로 乙의 동의가 있는 것처럼 믿게 하여 자기 소유 건물을 丙에게 매각하는 계약을 체결하였다. 이에 관한 설명으로 옳지 않은 것을 모두 고른 것은? (다툼이 있으면 판례에 따름)
\begin{framedtext}
ㄱ. 乙은 丙을 상대로 계약을 취소할 수 있다. \par
ㄴ. 丙은 甲을 상대로 계약의 추인여부에 대한 확답을 촉구할 수 있다. \par
ㄷ. 계약 당시 甲이 제한능력자임을 丙이 알았더라도 그 추인이 있기 전까지 丙은 乙을 상대로 자기의 의사표시를 철회할 수 있다. \par
\end{framedtext}
\begin{choices}
  \item ㄱ
  \item ㄷ
  \item ㄱ, ㄴ
  \item ㄴ, ㄷ
  \item ㄱ, ㄴ, ㄷ
\end{choices}
\end{question}
\answer{4}
